\uuid{7tRu}
\titre{Module/argument}
\auteur{bodin ; hanani ; mzari}
\organisation{exo7}
\licence{CC BY-NC-SA 4.0}
\datecreate{2020-01-20}
\niveau{L1}
\module{Algèbre}
\chapitre{Nombres complexes}
\sousChapitre{Racine carrée, équation du second degré}
\difficulte{3}
\theme{Équations, module, argument}
\exo7qcmid{lille-1.nombres-complexes.028}

\contenu{
\begin{question}
\qclassification{104.02,104.03,104.04}
\qtags{motcle=module/argument}
Soit $E$ l'ensemble des points $M$ d'affixe $z$ tels que : $|\frac{z-1}{1+iz}|=\sqrt 2$. Quelles sont les assertions vraies ?
\begin{answers}
    \bad{$E$ est une droite.}

    \good{$E$ est un cercle.}

    \bad{$E=\emptyset$}

    \good{$E$ est le  cercle de rayon $2$ et de centre le point d'affixe $-1+2i$.}

\end{answers}
\begin{explanations}
Soit $z \neq i$. On a :  $|\frac{z-1}{1+iz}|=\sqrt 2 \Leftrightarrow |z-1|^2=2|1+iz|^2 \Leftrightarrow
(z-1)(\overline{z}-1)=2 (1+iz)(1-i\overline{z})$. On développe cette dernière égalité.
\end{explanations}
\end{question}
}
